\section{\texorpdfstring{$D_\pm$}{D±} 型から得られる第三の全射影成分}\label{sec:dplus-full}
前節の共通正規形で $D_+$ の係数を取り、$p=x^2+2xz+3z^2$、
$P_D=x_0^2J+Q_D$ とする。
補助空間の成分であることに加え、全150係数でも次が成立する。

\begin{theorem}\label{thm:dplus-full}
$[P_D]$ の $\PGL_5$ 軌道閉包は、既出の54成分とは異なる23次元の全 $\cP$ 成分である。
代表点で
\[
 \widehat\OO_{\cP,[P_D]}
       \simeq\CC[[t_1,\ldots,t_{23},\varepsilon]]/(\varepsilon^3).
\]
従って $\#\Irr(\cP)\ge55$ である。
\end{theorem}
\begin{proof}
補助空間で得た二次 jet $Q_D+\varepsilon R+\varepsilon^2T$ を用いる。
$[J,R]=[J,T]=0$ なので斉次化にも同じ二次 jet が使える。
$\QQ(\sqrt3)$ 上で全350方程式・150変数の微分を作ると、
階数は101、垂直方向の階数は24、実際の群作用・垂直方向・スカラーの合併階数は48である。
$R$ を追加すると49となり全接空間を張る。
二次補正の全曖昧さを追加しても微分の階数は101のままで、第三次係数を追加すると102となる。
全計算は別の二つの多項式実装で一致した。

実際の $\PGL_5$ 軌道の次元は23であり、補正を群作用で運べるので、
補題\ref{lem:universal-two-jet-v7}が適用できる。
これで全成分性、次元、局所環が従う。残る相互区別は以下の退化因子で行う。
\end{proof}

\subsection{三つの cusp 平面を持つ既約四次錐}
$t=x_0$、$r=\sqrt3$ と書く。退化因子は $D_D=tF_D$、
\begin{align*}
 F_D&=t^4+t^2xz+p_D(x,z),\\
 p_D&=-\frac r{864}x^4+\frac1{36}x^3z+\frac18x^2z^2
           +\frac16xz^3+\frac r{24}z^4.
\end{align*}
ここで
\begin{align*}
 p_D&=-\frac r{864}\bigl(x-(9+5r)z\bigr)
                         \bigl(x+(3-r)z\bigr)^3,\\
 \Delta=(xz)^2-4p_D&=\frac r{216}\bigl(x-(3+r)z\bigr)^3
                              \bigl(x+(9-5r)z\bigr).
\end{align*}
$p_D,\Delta$ はいずれも異なる一次因子の重複度が $(3,1)$ なので、$\CC(x,z)$ で平方でない。
偶四次式 $t^4+st^2+p$ が可約なら、その二次因子分解を
$(t^2+at+b)(t^2-at+c)$ と書ける。
一次因子がある場合も符号を変えた根があるので、この場合に含まれる。
$t$ の係数 $a(c-b)=0$ から、$a=0$ なら $s^2-4p=(b-c)^2$、
$b=c$ なら $p=b^2$ である。従って $F_D$ は $\CC(x,z)[t]$ で既約、
Gauss の補題により $\CC[t,x,z]$ でも既約である。

平面曲線 $F_D=0$ の特異点は正確に
\[
 [0:r-3:1],\qquad [\gamma:3+r:1],\quad[-\gamma:3+r:1],
 \qquad \gamma^2=-\frac{3+r}{2}
\]
の三つである。
実際、$\partial_tF_D=2t(2t^2+xz)$ から、$t=0$ では $p_D$ の唯一の重根を取り、
$t\ne0$ では $t^2=-xz/2$ と $\Delta=\Delta_x=\Delta_z=0$ に帰着する。
後者では $\Delta$ の唯一の重根から二点が生じる。

第一点では $t^2$ の係数が非零で、もう一方向の最初の項は三次である。
他の二点では
\[
 4F_D=(2t^2+xz)^2-\Delta
\]
を使えば、$2t^2+xz$ と $x-(3+r)z$ が局所座標になり、同じく $A_2$ cusp と分かる。
$\PP^4$ 内の特異 locus は、それら三点と頂点直線を結んだ三つの $\PP^2$ である。
第一の平面 $\Pi_0=\{t=x+(3-r)z=0\}$ に沿い $D_D\in I_{\Pi_0}^3$、
他の二平面に沿う一般重複度は2である。

\subsection{相異なる55成分という下界}
ある平面に沿う五次因子の3次以上の消滅も閉 incidence 条件であり、
SNC 因子と両立しない。従って本成分は MPS の40成分とは異なる。
一般 rank と次元で、六つの rank 2 成分、二つの22次元成分、二つの28次元成分を除く。
$H+F_4$ という分解は五重超平面の成分と異なる。
v6 の残る23次元成分の四次部分は本質五変数である。
今回の $M,A_0$ は特異平面が一つで、その一般四次重複度は3であるのに対し、
$D_\pm$ は特異平面が三つで、すべて一般四次重複度2である。
いずれも開軌道を持つので、これらの異なる一般不変量を持つ軌道閉包が同一成分になることはない。
以上より
\[
 \#\Irr(\cP)\ \ge\ 52+1_M+1_{A_0}+1_{D_\pm}=55.
\]
